% --- Archon physics formalization source begin ---
% archon:physics
% archon:covers PhyXMiniProblems/problem_phyx_mini_0813.lean
% archon:source-report reports/phyx_mini/problem_phyx_mini_0813.source.json

\chapter{Physics problem phyx\_mini\_0813}
\label{ch:PhyXMiniProblems_problem_phyx_mini_0813}

\paragraph{Problem source.}
\#\# Physical scenario

A glider with mass \textbackslash{}( m = 0.200 \textbackslash{}) kg sits on a frictionless, horizontal air track, connected to a spring with force constant \textbackslash{}( k = 5.00 \textbackslash{}) N/m. You pull on the glider, stretching the spring 0.100 m, and release it from rest as shown in figure. The glider moves back toward its equilibrium position (\textbackslash{}( x = 0 \textbackslash{})).

\#\# Question

What is its \textbackslash{}( x \textbackslash{})-velocity when \textbackslash{}( x = 0.080 \textbackslash{}) m?

\#\# Answer choices

- A: A: \textbackslash{}( -0.70 \textbackslash{}text\{ m/s\} \textbackslash{})

- B: B: \textbackslash{}( -0.30 \textbackslash{}text\{ m/s\} \textbackslash{})

- C: C: \textbackslash{}( -1.20 \textbackslash{}text\{ m/s\} \textbackslash{})

- D: D: \textbackslash{}( -2.30 \textbackslash{}text\{ m/s\} \textbackslash{})

\#\# Figure caption (auxiliary; use the image as primary evidence)

The image illustrates a spring-mass system with the following details:

1. **Spring Details**:
   - The spring constant is given as \textbackslash{}( k = 5.00 \textbackslash{}, \textbackslash{}text\{N/m\} \textbackslash{}).
   - The spring is shown in a relaxed position labeled "Spring relaxed."

2. **Position and Displacement**:
   - The relaxed position of the spring is marked at \textbackslash{}( x = 0 \textbackslash{}).
   - A block is positioned at \textbackslash{}( x\_1 = 0.100 \textbackslash{}, \textbackslash{}text\{m\} \textbackslash{}) from the relaxed position.

3. **Block Details**:
   - The block has a mass labeled as \textbackslash{}( m = 0.200 \textbackslash{}, \textbackslash{}text\{kg\} \textbackslash{}).
   - The initial velocity of the block is \textbackslash{}( v\_\{1x\} = 0 \textbackslash{}).

4. **Axes and Labels**:
   - The horizontal axis is denoted as the \textbackslash{}( x \textbackslash{})-axis.
   - "Point 1" is labeled at the beginning of the spring.

The visual emphasizes the displacement of the block from the equilibrium position of the spring, with given parameters related to force and motion.

\#\# Dataset metadata

Work and Energy; Physical Model Grounding Reasoning, Multi-Formula Reasoning

\paragraph{Recorded answer/context.}
B: B: \textbackslash{}( -0.30 \textbackslash{}text\{ m/s\} \textbackslash{})

\paragraph{Figure/image path.}
/root/proposal\_for\_physic/science-mango/phyx\_mini\_run/phyx\_data/test\_image/813.png

\paragraph{Formalization target.}
create a compiling Lean file with sorry bodies at `PhyXMiniProblems/problem\_phyx\_mini\_0813.lean`. The Lean declarations must preserve the physical quantities, dimensions or dimensional roles, figure labels, governing-law hypotheses, and final relation expressed by this problem.
Use Mathlib/Physlib names found through LeanExplore where available. If a physics API is missing, introduce faithful local abstractions rather than scalar placeholder aliases.

\begin{theorem}[Physics formalization target]
\label{thm:physics:phyx_mini_0813:target}
\lean{PhyXMiniProblems.ProblemPhyXMini0813.problem_phyx_mini_0813}
\uses{def:physics:phyx-mini-0813:phyxminiproblems-problemphyxmini0813-velocityinmeterspersecond, def:physics:phyx-mini-0813:phyxminiproblems-problemphyxmini0813-gliderspringsetup, def:physics:phyx-mini-0813:phyxminiproblems-problemphyxmini0813-matchesproblemandprimaryfigure, def:physics:phyx-mini-0813:phyxminiproblems-problemphyxmini0813-hasphysicalparameters, def:physics:phyx-mini-0813:phyxminiproblems-problemphyxmini0813-satisfiesfrictionlessharmonicmotion, def:physics:phyx-mini-0813:phyxminiproblems-problemphyxmini0813-recordeddatasetanswerchoice, def:physics:phyx-mini-0813:phyxminiproblems-problemphyxmini0813-answerchoicereportsqueriedvelocity}
The assigned autoformalize agent should translate this physics problem into Lean declarations in the covered file, with theorem and lemma proof bodies written as `by sorry`.
\end{theorem}
\begin{proof}
This is an autoformalization task, not a proof task. Produce statements that can later be proved by the physics prover without weakening the physical model.
\end{proof}

% --- Archon declaration topology begin ---
\section{Declaration topology}
\begin{definition}[Mass Quantity]
  \label{def:physics:phyx-mini-0813:phyxminiproblems-problemphyxmini0813-massquantity}
  \lean{PhyXMiniProblems.ProblemPhyXMini0813.MassQuantity}
  A nonnegative, unit-independent physical mass
\end{definition}
\begin{proof}
  This is the declared model definition of Mass Quantity.
\end{proof}
\begin{definition}[Signed Position Quantity]
  \label{def:physics:phyx-mini-0813:phyxminiproblems-problemphyxmini0813-signedpositionquantity}
  \lean{PhyXMiniProblems.ProblemPhyXMini0813.SignedPositionQuantity}
  A signed position along the horizontal x-axis
\end{definition}
\begin{proof}
  This is the declared model definition of Signed Position Quantity.
\end{proof}
\begin{definition}[Signed Velocity Quantity]
  \label{def:physics:phyx-mini-0813:phyxminiproblems-problemphyxmini0813-signedvelocityquantity}
  \lean{PhyXMiniProblems.ProblemPhyXMini0813.SignedVelocityQuantity}
  A signed horizontal velocity, with physical dimension L T⁻¹
\end{definition}
\begin{proof}
  This is the declared model definition of Signed Velocity Quantity.
\end{proof}
\begin{definition}[Spring Stiffness Quantity]
  \label{def:physics:phyx-mini-0813:phyxminiproblems-problemphyxmini0813-springstiffnessquantity}
  \lean{PhyXMiniProblems.ProblemPhyXMini0813.SpringStiffnessQuantity}
  A nonnegative spring stiffness, with dimension M T⁻² (N/m)
\end{definition}
\begin{proof}
  This is the declared model definition of Spring Stiffness Quantity.
\end{proof}
\begin{definition}[mass Readout]
  \label{def:physics:phyx-mini-0813:phyxminiproblems-problemphyxmini0813-massreadout}
  \lean{PhyXMiniProblems.ProblemPhyXMini0813.massReadout}
  \uses{def:physics:phyx-mini-0813:phyxminiproblems-problemphyxmini0813-massquantity}
  Read a physical mass in a selected mass unit
\end{definition}
\begin{proof}
  This is the declared model definition of mass Readout.
\end{proof}
\begin{definition}[position Readout]
  \label{def:physics:phyx-mini-0813:phyxminiproblems-problemphyxmini0813-positionreadout}
  \lean{PhyXMiniProblems.ProblemPhyXMini0813.positionReadout}
  \uses{def:physics:phyx-mini-0813:phyxminiproblems-problemphyxmini0813-signedpositionquantity}
  Read a signed x-position in a selected length unit
\end{definition}
\begin{proof}
  This is the declared model definition of position Readout.
\end{proof}
\begin{definition}[velocity Readout]
  \label{def:physics:phyx-mini-0813:phyxminiproblems-problemphyxmini0813-velocityreadout}
  \lean{PhyXMiniProblems.ProblemPhyXMini0813.velocityReadout}
  \uses{def:physics:phyx-mini-0813:phyxminiproblems-problemphyxmini0813-signedvelocityquantity}
  Read a signed x-velocity in coherent selected length and time units
\end{definition}
\begin{proof}
  This is the declared model definition of velocity Readout.
\end{proof}
\begin{definition}[spring Stiffness Readout]
  \label{def:physics:phyx-mini-0813:phyxminiproblems-problemphyxmini0813-springstiffnessreadout}
  \lean{PhyXMiniProblems.ProblemPhyXMini0813.springStiffnessReadout}
  \uses{def:physics:phyx-mini-0813:phyxminiproblems-problemphyxmini0813-springstiffnessquantity}
  Read spring stiffness in coherent selected mass and time units
\end{definition}
\begin{proof}
  This is the declared model definition of spring Stiffness Readout.
\end{proof}
\begin{definition}[mass In Kilograms]
  \label{def:physics:phyx-mini-0813:phyxminiproblems-problemphyxmini0813-massinkilograms}
  \lean{PhyXMiniProblems.ProblemPhyXMini0813.massInKilograms}
  \uses{def:physics:phyx-mini-0813:phyxminiproblems-problemphyxmini0813-massquantity, def:physics:phyx-mini-0813:phyxminiproblems-problemphyxmini0813-massreadout}
  Kilogram readout of the glider mass
\end{definition}
\begin{proof}
  This is the declared model definition of mass In Kilograms.
\end{proof}
\begin{definition}[position In Meters]
  \label{def:physics:phyx-mini-0813:phyxminiproblems-problemphyxmini0813-positioninmeters}
  \lean{PhyXMiniProblems.ProblemPhyXMini0813.positionInMeters}
  \uses{def:physics:phyx-mini-0813:phyxminiproblems-problemphyxmini0813-signedpositionquantity, def:physics:phyx-mini-0813:phyxminiproblems-problemphyxmini0813-positionreadout}
  Metre readout of a signed x-position
\end{definition}
\begin{proof}
  This is the declared model definition of position In Meters.
\end{proof}
\begin{definition}[velocity In Meters Per Second]
  \label{def:physics:phyx-mini-0813:phyxminiproblems-problemphyxmini0813-velocityinmeterspersecond}
  \lean{PhyXMiniProblems.ProblemPhyXMini0813.velocityInMetersPerSecond}
  \uses{def:physics:phyx-mini-0813:phyxminiproblems-problemphyxmini0813-signedvelocityquantity, def:physics:phyx-mini-0813:phyxminiproblems-problemphyxmini0813-velocityreadout}
  Metres-per-second readout, positive in the figure's positive x-direction
\end{definition}
\begin{proof}
  This is the declared model definition of velocity In Meters Per Second.
\end{proof}
\begin{definition}[spring Stiffness In Newtons Per Meter]
  \label{def:physics:phyx-mini-0813:phyxminiproblems-problemphyxmini0813-springstiffnessinnewtonspermeter}
  \lean{PhyXMiniProblems.ProblemPhyXMini0813.springStiffnessInNewtonsPerMeter}
  \uses{def:physics:phyx-mini-0813:phyxminiproblems-problemphyxmini0813-springstiffnessquantity, def:physics:phyx-mini-0813:phyxminiproblems-problemphyxmini0813-springstiffnessreadout}
  Newton-per-metre readout of the spring stiffness
\end{definition}
\begin{proof}
  This is the declared model definition of spring Stiffness In Newtons Per Meter.
\end{proof}
\begin{definition}[Process Instant]
  \label{def:physics:phyx-mini-0813:phyxminiproblems-problemphyxmini0813-processinstant}
  \lean{PhyXMiniProblems.ProblemPhyXMini0813.ProcessInstant}
  The release state and the later state asked about in the problem
\end{definition}
\begin{proof}
  This is the declared model definition of Process Instant.
\end{proof}
\begin{definition}[Figure Feature]
  \label{def:physics:phyx-mini-0813:phyxminiproblems-problemphyxmini0813-figurefeature}
  \lean{PhyXMiniProblems.ProblemPhyXMini0813.FigureFeature}
  Features visibly drawn in the supplied image
\end{definition}
\begin{proof}
  This is the declared model definition of Figure Feature.
\end{proof}
\begin{definition}[Figure Label]
  \label{def:physics:phyx-mini-0813:phyxminiproblems-problemphyxmini0813-figurelabel}
  \lean{PhyXMiniProblems.ProblemPhyXMini0813.FigureLabel}
  Text and mathematical labels printed in the supplied image
\end{definition}
\begin{proof}
  This is the declared model definition of Figure Label.
\end{proof}
\begin{definition}[Figure Label Role]
  \label{def:physics:phyx-mini-0813:phyxminiproblems-problemphyxmini0813-figurelabelrole}
  \lean{PhyXMiniProblems.ProblemPhyXMini0813.FigureLabelRole}
  Physical meanings of the labels in image 813.png
\end{definition}
\begin{proof}
  This is the declared model definition of Figure Label Role.
\end{proof}
\begin{definition}[Horizontal Direction]
  \label{def:physics:phyx-mini-0813:phyxminiproblems-problemphyxmini0813-horizontaldirection}
  \lean{PhyXMiniProblems.ProblemPhyXMini0813.HorizontalDirection}
  Horizontal orientations used for the track and x-axis
\end{definition}
\begin{proof}
  This is the declared model definition of Horizontal Direction.
\end{proof}
\begin{definition}[Track Condition]
  \label{def:physics:phyx-mini-0813:phyxminiproblems-problemphyxmini0813-trackcondition}
  \lean{PhyXMiniProblems.ProblemPhyXMini0813.TrackCondition}
  The idealized support condition stated in the prose
\end{definition}
\begin{proof}
  This is the declared model definition of Track Condition.
\end{proof}
\begin{definition}[Glider Spring Figure]
  \label{def:physics:phyx-mini-0813:phyxminiproblems-problemphyxmini0813-gliderspringfigure}
  \lean{PhyXMiniProblems.ProblemPhyXMini0813.GliderSpringFigure}
  \uses{def:physics:phyx-mini-0813:phyxminiproblems-problemphyxmini0813-figurefeature, def:physics:phyx-mini-0813:phyxminiproblems-problemphyxmini0813-figurelabel, def:physics:phyx-mini-0813:phyxminiproblems-problemphyxmini0813-figurelabelrole, def:physics:phyx-mini-0813:phyxminiproblems-problemphyxmini0813-horizontaldirection}
  Qualitative transcription of the primary image
\end{definition}
\begin{proof}
  This is the declared model definition of Glider Spring Figure.
\end{proof}
\begin{definition}[Glider Spring Setup]
  \label{def:physics:phyx-mini-0813:phyxminiproblems-problemphyxmini0813-gliderspringsetup}
  \lean{PhyXMiniProblems.ProblemPhyXMini0813.GliderSpringSetup}
  \uses{def:physics:phyx-mini-0813:phyxminiproblems-problemphyxmini0813-massquantity, def:physics:phyx-mini-0813:phyxminiproblems-problemphyxmini0813-signedpositionquantity, def:physics:phyx-mini-0813:phyxminiproblems-problemphyxmini0813-signedvelocityquantity, def:physics:phyx-mini-0813:phyxminiproblems-problemphyxmini0813-springstiffnessquantity, def:physics:phyx-mini-0813:phyxminiproblems-problemphyxmini0813-processinstant, def:physics:phyx-mini-0813:phyxminiproblems-problemphyxmini0813-trackcondition, def:physics:phyx-mini-0813:phyxminiproblems-problemphyxmini0813-gliderspringfigure}
  Independent physical quantities and the associated one-dimensional Physlib trajectory. The velocity at the queried position is an unconstrained physical quantity here; it is not defined from an answer choice
\end{definition}
\begin{proof}
  This is the declared model definition of Glider Spring Setup.
\end{proof}
\begin{definition}[Matches Problem And Primary Figure]
  \label{def:physics:phyx-mini-0813:phyxminiproblems-problemphyxmini0813-matchesproblemandprimaryfigure}
  \lean{PhyXMiniProblems.ProblemPhyXMini0813.MatchesProblemAndPrimaryFigure}
  \uses{def:physics:phyx-mini-0813:phyxminiproblems-problemphyxmini0813-massinkilograms, def:physics:phyx-mini-0813:phyxminiproblems-problemphyxmini0813-positioninmeters, def:physics:phyx-mini-0813:phyxminiproblems-problemphyxmini0813-velocityinmeterspersecond, def:physics:phyx-mini-0813:phyxminiproblems-problemphyxmini0813-springstiffnessinnewtonspermeter, def:physics:phyx-mini-0813:phyxminiproblems-problemphyxmini0813-gliderspringsetup}
  Numerical and qualitative data stated by the source and shown in image 813.png. The strict negative velocity records only the stated direction "back toward equilibrium" from the positive-x side; it does not prescribe the requested speed
\end{definition}
\begin{proof}
  This is the declared model definition of Matches Problem And Primary Figure.
\end{proof}
\begin{definition}[Has Physical Parameters]
  \label{def:physics:phyx-mini-0813:phyxminiproblems-problemphyxmini0813-hasphysicalparameters}
  \lean{PhyXMiniProblems.ProblemPhyXMini0813.HasPhysicalParameters}
  \uses{def:physics:phyx-mini-0813:phyxminiproblems-problemphyxmini0813-massinkilograms, def:physics:phyx-mini-0813:phyxminiproblems-problemphyxmini0813-springstiffnessinnewtonspermeter, def:physics:phyx-mini-0813:phyxminiproblems-problemphyxmini0813-gliderspringsetup}
  Positivity conditions selecting a physical oscillator
\end{definition}
\begin{proof}
  This is the declared model definition of Has Physical Parameters.
\end{proof}
\begin{definition}[Satisfies Frictionless Harmonic Motion]
  \label{def:physics:phyx-mini-0813:phyxminiproblems-problemphyxmini0813-satisfiesfrictionlessharmonicmotion}
  \lean{PhyXMiniProblems.ProblemPhyXMini0813.SatisfiesFrictionlessHarmonicMotion}
  \uses{def:physics:phyx-mini-0813:phyxminiproblems-problemphyxmini0813-massinkilograms, def:physics:phyx-mini-0813:phyxminiproblems-problemphyxmini0813-positioninmeters, def:physics:phyx-mini-0813:phyxminiproblems-problemphyxmini0813-velocityinmeterspersecond, def:physics:phyx-mini-0813:phyxminiproblems-problemphyxmini0813-springstiffnessinnewtonspermeter, def:physics:phyx-mini-0813:phyxminiproblems-problemphyxmini0813-gliderspringsetup}
  The physical readouts realize a one-dimensional Physlib harmonic oscillator. Its coordinate is displacement from the relaxed-spring equilibrium, its time derivative is the signed x-velocity, and its total mechanical energy is equal at release and at the queried position because the horizontal track is frictionless. The energy equality is a governing conservation law. It does not contain a numerical value for the queried velocity
\end{definition}
\begin{proof}
  This is the declared model definition of Satisfies Frictionless Harmonic Motion.
\end{proof}
\begin{lemma}[queried Velocity Squared eq nine Hundredths]
  \label{lem:physics:phyx-mini-0813:phyxminiproblems-problemphyxmini0813-queriedvelocitysquared-eq-ninehundredths}
  \lean{PhyXMiniProblems.ProblemPhyXMini0813.queriedVelocitySquared_eq_nineHundredths}
  \uses{def:physics:phyx-mini-0813:phyxminiproblems-problemphyxmini0813-velocityinmeterspersecond, def:physics:phyx-mini-0813:phyxminiproblems-problemphyxmini0813-gliderspringsetup, def:physics:phyx-mini-0813:phyxminiproblems-problemphyxmini0813-matchesproblemandprimaryfigure, def:physics:phyx-mini-0813:phyxminiproblems-problemphyxmini0813-satisfiesfrictionlessharmonicmotion}
  Conservation of 1/2 m v² + 1/2 k x² gives the squared queried speed (0.30 m/s)² = 9/100 (m/s)². This magnitude is derived rather than assumed
\end{lemma}
\begin{proof}
  The conclusion follows from the governing relations and figure readouts named in the statement of queried Velocity Squared eq nine Hundredths.
\end{proof}
\begin{definition}[Answer Choice]
  \label{def:physics:phyx-mini-0813:phyxminiproblems-problemphyxmini0813-answerchoice}
  \lean{PhyXMiniProblems.ProblemPhyXMini0813.AnswerChoice}
  Labels of the four displayed answer choices
\end{definition}
\begin{proof}
  This is the declared model definition of Answer Choice.
\end{proof}
\begin{definition}[velocity In Meters Per Second]
  \label{def:physics:phyx-mini-0813:phyxminiproblems-problemphyxmini0813-answerchoice-velocityinmeterspersecond}
  \lean{PhyXMiniProblems.ProblemPhyXMini0813.AnswerChoice.velocityInMetersPerSecond}
  \uses{def:physics:phyx-mini-0813:phyxminiproblems-problemphyxmini0813-velocityinmeterspersecond, def:physics:phyx-mini-0813:phyxminiproblems-problemphyxmini0813-answerchoice}
  Signed x-velocity in metres per second printed beside each choice
\end{definition}
\begin{proof}
  This is the declared model definition of velocity In Meters Per Second.
\end{proof}
\begin{definition}[recorded Dataset Answer Choice]
  \label{def:physics:phyx-mini-0813:phyxminiproblems-problemphyxmini0813-recordeddatasetanswerchoice}
  \lean{PhyXMiniProblems.ProblemPhyXMini0813.recordedDatasetAnswerChoice}
  \uses{def:physics:phyx-mini-0813:phyxminiproblems-problemphyxmini0813-answerchoice}
  The answer label recorded in the supplied dataset metadata
\end{definition}
\begin{proof}
  This is the declared model definition of recorded Dataset Answer Choice.
\end{proof}
\begin{definition}[Answer Choice Reports Queried Velocity]
  \label{def:physics:phyx-mini-0813:phyxminiproblems-problemphyxmini0813-answerchoicereportsqueriedvelocity}
  \lean{PhyXMiniProblems.ProblemPhyXMini0813.AnswerChoiceReportsQueriedVelocity}
  \uses{def:physics:phyx-mini-0813:phyxminiproblems-problemphyxmini0813-velocityinmeterspersecond, def:physics:phyx-mini-0813:phyxminiproblems-problemphyxmini0813-gliderspringsetup, def:physics:phyx-mini-0813:phyxminiproblems-problemphyxmini0813-answerchoice}
  A displayed choice reports the glider's queried signed x-velocity
\end{definition}
\begin{proof}
  This is the declared model definition of Answer Choice Reports Queried Velocity.
\end{proof}
% --- Archon declaration topology end ---
% --- Archon physics formalization source end ---
